\section{总结与延伸}

本节把整期访谈收束成五个结论。第一，具身智能最初不是传统机器人产业的自然延伸，而是计算机视觉从被动识别走向主动行动的学术迁移。第二，语言不是智能本质，而是人类智能的一次跃变；机器人智能必须回到环境反应和物理反馈。第三，VLM/VLA 的瓶颈很大程度来自数据覆盖不足，尤其是 action data 稀缺。第四，软件、硬件、数据和场景必须螺旋上升，不能单独押注通用本体或通用大脑。第五，产业最终必须用生产力、出货和真实展示证明自己。

\begin{importantbox}{把 EP106 放进张小珺机器人线索}
EP106 提供学术史和行业伦理，EP109 提供仿真/合成数据基础设施，EP121 提供 DeepMind 机器人模型视角，EP132 提供整机和供应链案例。四期合在一起，构成“概念起源--数据生产--模型路线--产业落地”的机器人学习路径。
\end{importantbox}

\subsection{关键 takeaways}

\begin{enumerate}
\item 具身智能的核心不是机器人外形，而是感知、行动和环境反馈的闭环。
\item 语言增强智能，但物理世界智能必须通过身体和 action data 接地。
\item 真实数据很贵，合成数据和出货回流是机器人数据 recipe 的关键组成。
\item 早期公司要避免长期漂浮和算不过账，先在应用场景内做泛化。
\item 行业信用来自真实展示和生产力验证，遥操作、融资和 demo 不能替代真实进展。
\end{enumerate}

\subsection{拓展阅读}

\begin{itemize}
\item 想理解仿真和合成数据基础设施，可对照 EP109 谢晨访谈。
\item 想理解机器人基座模型、跨本体和世界模型，可对照 EP121 DeepMind 谭捷访谈。
\item 想理解整机、供应链和 Data Recipe，可对照 EP132 星海图高继扬访谈。
\end{itemize}

\end{document}
